\section{Polynomials}

\subsection{Definition}

\begin{definition}
\label{def.fps.pol}
\lean{FPS.IsPolynomial, FPS.polynomialSubalgebra}
\leantarget
\leanok
\textbf{(a)} An FPS $a\in K\left[  \left[  x\right]
\right]  $ is said to be a \emph{polynomial} if all but finitely many
$n\in\mathbb{N}$ satisfy $\left[  x^{n}\right]  a=0$ (that is, if all but
finitely many coefficients of $a$ are $0$). \medskip

\textbf{(b)} We let $K\left[  x\right]  $ be the set of all polynomials $a\in
K\left[  \left[  x\right]  \right]  $. This set $K\left[  x\right]  $ is a
subring of $K\left[  \left[  x\right]  \right]  $ (according to Theorem
\ref{thm.fps.pol.ring} below), and is called the \emph{univariate polynomial
ring} over $K$.
\end{definition}

\begin{lemma}
\label{lem.fps.pol.degree-bound}
\lean{FPS.isPolynomial_iff_exists_degree_bound}
\leanhelper
\leanok
An FPS $f\in K\left[\left[x\right]\right]$ is a polynomial if and only if
there exists an $N\in\mathbb{N}$ such that $\left[x^{n}\right]f=0$ for all
$n\geq N$.
\end{lemma}

\begin{proof}
\leanok
$\Longrightarrow$: If the set of nonzero-coefficient indices is finite, it is
bounded above by some $N$, so all coefficients beyond $N$ vanish.

$\Longleftarrow$: If all coefficients beyond $N$ vanish, then the set of
nonzero-coefficient indices is contained in $\{0,1,\ldots,N-1\}$, which is
finite.
\end{proof}

\begin{lemma}
\label{lem.fps.pol.exists-polynomial}
\lean{FPS.isPolynomial_iff_exists_polynomial}
\leanhelper
\leanok
An FPS $f\in K\left[\left[x\right]\right]$ is a polynomial (in the sense of
Definition~\ref{def.fps.pol}) if and only if $f$ equals the image of some
element of $K[x]$ under the canonical embedding $K[x]\hookrightarrow
K\left[\left[x\right]\right]$.
\end{lemma}

\begin{proof}
\leanok
$\Longrightarrow$: Use the degree bound from
Lemma~\ref{lem.fps.pol.degree-bound} and truncation.

$\Longleftarrow$: The image of any polynomial has finite support
(Lemma~\ref{lem.fps.pol.of-polynomial}).
\end{proof}

\begin{lemma}
\label{lem.fps.pol.of-polynomial}
\lean{FPS.isPolynomial_of_polynomial}
\leanhelper
\leanok
The image of any polynomial $p\in K[X]$ under the canonical embedding
$K[X]\hookrightarrow K\left[\left[x\right]\right]$ is a polynomial in the FPS
sense.
\end{lemma}

\begin{proof}
\leanok
The set of indices with nonzero coefficients is contained in the support of $p$,
which is a finite set.
\end{proof}

\subsubsection{Converting between polynomial FPSs and polynomials}

\begin{definition}
\label{def.fps.pol.toPolynomial}
\lean{FPS.toPolynomial}
\leanhelper
\leanok
Let $f\in K\left[\left[x\right]\right]$ be a polynomial FPS (in the sense of
Definition~\ref{def.fps.pol}). We define $\operatorname{toPolynomial}(f)\in K[X]$ as
the truncation of $f$ at the degree bound given by
Lemma~\ref{lem.fps.pol.degree-bound}.
\end{definition}

\begin{lemma}
\label{lem.fps.pol.coe-toPolynomial}
\lean{FPS.coe_toPolynomial}
\leanhelper
\leanok
Let $f\in K\left[\left[x\right]\right]$ be a polynomial FPS. Then the
coercion of $\operatorname{toPolynomial}(f)$ back into
$K\left[\left[x\right]\right]$ equals $f$:
\[
\iota\bigl(\operatorname{toPolynomial}(f)\bigr) = f,
\]
where $\iota\colon K[X]\hookrightarrow K\left[\left[x\right]\right]$ is the
canonical embedding.
\end{lemma}

\begin{proof}
By the degree bound, all coefficients of $f$ beyond the bound are zero, so
truncation and coercion recover every coefficient of $f$.
\end{proof}

\begin{lemma}
\label{lem.fps.pol.toPolynomial-coe}
\lean{FPS.toPolynomial_coe}
\leanhelper
\leanok
For any polynomial $p\in K[X]$, we have
$\operatorname{toPolynomial}(\iota(p))=p$, where $\iota$ is the canonical embedding.
\end{lemma}

\begin{proof}
\leanok
By Lemma~\ref{lem.fps.pol.coe-toPolynomial}, both sides have the same image
under $\iota$, and $\iota$ is injective.
\end{proof}

\subsection{Polynomials form a subring of power series}

\begin{lemma}
\label{lem.fps.pol.zero}
\lean{FPS.isPolynomial_zero}
\leanhelper
\leanok
The zero power series $\underline{0}$ is a polynomial.
\end{lemma}

\begin{proof}
\leanok
All coefficients of $\underline{0}$ are $0$, so the set of nonzero-coefficient
indices is empty, hence finite.
\end{proof}

\begin{lemma}
\label{lem.fps.pol.one}
\lean{FPS.isPolynomial_one}
\leanhelper
\leanok
The unity $\underline{1}$ is a polynomial.
\end{lemma}

\begin{proof}
\leanok
The unity $\underline{1}$ equals the image of $1\in K[X]$ under the canonical
embedding.
\end{proof}

\begin{lemma}
\label{lem.fps.pol.add}
\lean{FPS.isPolynomial_add}
\leanhelper
\leanok
The sum of two polynomial power series is a polynomial.
\end{lemma}

\begin{proof}
\leanok
If $f$ and $g$ are polynomials, then the set of indices where $(f+g)$ has a
nonzero coefficient is contained in the union of the corresponding sets for
$f$ and $g$. A finite union of finite sets is finite.
\end{proof}

\begin{lemma}
\label{lem.fps.pol.neg}
\lean{FPS.isPolynomial_neg}
\leanhelper
\leanok
The negation of a polynomial power series is a polynomial.
\end{lemma}

\begin{proof}
\leanok
Negation does not change which coefficients are nonzero.
\end{proof}

\begin{lemma}
\label{lem.fps.pol.sub}
\lean{FPS.isPolynomial_sub}
\leanhelper
\leanok
The difference of two polynomial power series is a polynomial.
\end{lemma}

\begin{proof}
\leanok
Write $f-g=f+(-g)$ and apply Lemma~\ref{lem.fps.pol.add} and
Lemma~\ref{lem.fps.pol.neg}.
\end{proof}

\begin{lemma}
\label{lem.fps.pol.mul}
\lean{FPS.isPolynomial_mul}
\leanhelper
\leanok
The product of two polynomial power series is a polynomial.
\end{lemma}

\begin{proof}
\leanok
Let $a,b\in K\left[  x\right]  $. Since $a$ is a polynomial, there exists a finite
subset $I$ of $\mathbb{N}$ such that
$\left[  x^{i}\right]  a=0$ for all
$i\in\mathbb{N}\setminus I$.
Similarly, there exists a finite subset $J$ of
$\mathbb{N}$ such that
$\left[  x^{j}\right]  b=0$ for all
$j\in\mathbb{N}\setminus J$.

Let $S=\left\{
i+j\ \mid\ i\in I\text{ and }j\in J\right\}$. This set $S$
is again finite (since $I$ and $J$ are finite), and we have
\[
\left[  x^{n}\right]  \left(  ab\right)  =\sum_{i=0}^{n}\left[  x^{i}\right]
a\cdot\left[  x^{n-i}\right]  b=0\text{ for all }n\in\mathbb{N}\setminus S,
\]
because for each $i\in\left\{  0,1,\ldots,n\right\}$, letting $j=n-i$, we cannot have both
$i\in I$ and $j\in J$ (otherwise $n=i+j\in S$, contradicting
$n\notin S$), so at least one of $\left[  x^{i}\right]  a$ and
$\left[  x^{j}\right]  b$ is $0$, making each summand $0$.
Thus, all but finitely many
$n\in\mathbb{N}$ satisfy $\left[  x^{n}\right]  \left(  ab\right)  =0$, so $ab\in K\left[  x\right]  $.
\end{proof}

\begin{lemma}
\label{lem.fps.pol.smul}
\lean{FPS.isPolynomial_smul}
\leanhelper
\leanok
Scalar multiplication of a polynomial power series by an element of $K$ is a
polynomial.
\end{lemma}

\begin{proof}
\leanok
If $\left[x^{n}\right]f=0$, then $\left[x^{n}\right](cf)=c\cdot
\left[x^{n}\right]f=0$. So the nonzero-coefficient set of $cf$ is contained
in that of $f$, which is finite.
\end{proof}

\begin{theorem}
\label{thm.fps.pol.ring}
\lean{FPS.polynomialSubalgebra, FPS.polynomialSubring, FPS.polynomialSubmodule}
\leantarget
\leanok
The set $K\left[  x\right]  $ is a subring of
$K\left[  \left[  x\right]  \right]  $ (that is, it is closed under addition,
subtraction and multiplication, and contains the zero $\underline{0}$ and the
unity $\underline{1}$) and is a $K$-submodule of $K\left[  \left[  x\right]
\right]  $ (that is, it is closed under addition and scaling by elements of
$K$).
\end{theorem}

\begin{proof}
Follows from the preceding lemmas: $K[x]$ contains $\underline{0}$
(Lemma~\ref{lem.fps.pol.zero}) and $\underline{1}$
(Lemma~\ref{lem.fps.pol.one}), is closed under addition
(Lemma~\ref{lem.fps.pol.add}), multiplication (Lemma~\ref{lem.fps.pol.mul}),
and scalar multiplication (Lemma~\ref{lem.fps.pol.smul}).
\end{proof}

\begin{lemma}
\label{lem.fps.pol.X}
\lean{FPS.isPolynomial_X}
\leanhelper
\leanok
The polynomial $x\in K\left[\left[x\right]\right]$ is a polynomial.
\end{lemma}

\begin{proof}
\leanok
The only nonzero coefficient of $x$ is the coefficient of $x^1$, which is $1$.
Thus the set of nonzero-coefficient indices is contained in $\{1\}$, which is
finite.
\end{proof}

\begin{lemma}
\label{lem.fps.pol.C}
\lean{FPS.isPolynomial_C}
\leanhelper
\leanok
For any $c\in K$, the constant FPS $\underline{c}$ is a polynomial.
\end{lemma}

\begin{proof}
\leanok
The only possibly nonzero coefficient of $\underline{c}$ is the coefficient of
$x^0$. Thus the set of nonzero-coefficient indices is contained in $\{0\}$,
which is finite.
\end{proof}

\subsection{Reminders on rings and $K$-algebras}

\begin{definition}
\label{def.alg.ring}
\lean{FPS.ring_add_assoc, FPS.ring_add_comm, FPS.ring_add_zero, FPS.ring_zero_add, FPS.ring_add_neg, FPS.ring_mul_assoc, FPS.ring_left_distrib, FPS.ring_right_distrib, FPS.ring_mul_one, FPS.ring_one_mul, FPS.ring_mul_zero, FPS.ring_zero_mul}
\leantarget
\leanok
The notion of a \emph{ring} (also known as a \emph{noncommutative ring}) is
defined in the exact same way as we defined the notion of a commutative ring in
Definition~\ref{def.alg.commring}, except that the ``Commutativity of
multiplication'' axiom is removed.
\end{definition}

\begin{definition}
\label{def.alg.Kalg}
\lean{FPS.kalg_add_comm, FPS.kalg_add_assoc, FPS.kalg_add_zero, FPS.kalg_mul_assoc, FPS.kalg_left_distrib, FPS.kalg_right_distrib, FPS.kalg_mul_one, FPS.kalg_one_mul, FPS.kalg_smul_assoc, FPS.kalg_smul_add, FPS.kalg_add_smul, FPS.kalg_one_smul, FPS.kalg_zero_smul, FPS.kalg_smul_zero, FPS.kalg_smul_mul_assoc, FPS.kalg_mul_smul_comm, FPS.kalg_smul_mul_eq_mul_smul}
\leantarget
\leanok
A $K$\emph{-algebra} is a set $A$ equipped with four maps%
\begin{align*}
\oplus &  :A\times A\rightarrow A,\\
\ominus &  :A\times A\rightarrow A,\\
\odot &  :A\times A\rightarrow A,\\
\rightharpoonup &  :K\times A\rightarrow A
\end{align*}
and two elements $\overrightarrow{0}\in A$ and $\overrightarrow{1}\in A$
satisfying the following properties:

\begin{enumerate}
\item The set $A$, equipped with the maps $\oplus$, $\ominus$ and $\odot$ and
the two elements $\overrightarrow{0}$ and $\overrightarrow{1}$, is a
(noncommutative) ring.

\item The set $A$, equipped with the maps $\oplus$, $\ominus$ and
$\rightharpoonup$ and the element $\overrightarrow{0}$, is a $K$-module.

\item We have%
\begin{equation}
\lambda\rightharpoonup\left(  a\odot b\right)  =\left(  \lambda\rightharpoonup
a\right)  \odot b=a\odot\left(  \lambda\rightharpoonup b\right)
\label{eq.def.alg.Kalg.scaleinv}%
\end{equation}
for all $\lambda\in K$ and $a,b\in A$.
\end{enumerate}

(Thus, in a nutshell, a $K$-algebra is a set $A$ that is simultaneously a ring
and a $K$-module, with the property that the ring $A$ and the $K$-module $A$
have the same addition, the same subtraction and the same zero, and satisfy
the additional compatibility property (\ref{eq.def.alg.Kalg.scaleinv}).)
\end{definition}

\subsection{Evaluation aka substitution into polynomials}

\begin{definition}
\label{def.pol.subs}
\lean{FPS.polyEval}
\leantarget
\leanok
Let $f\in K\left[  x\right]  $ be a polynomial. Let $A$ be
any $K$-algebra. Let $a\in A$ be any element. We then define an element
$f\left[  a\right]  \in A$ as follows:

Write $f$ in the form $f=\sum_{n\in\mathbb{N}}f_{n}x^{n}$ with $f_{0}%
,f_{1},f_{2},\ldots\in K$. (That is, $f_{n}=\left[  x^{n}\right]  f$ for each
$n\in\mathbb{N}$.) Then, set%
\[
f\left[  a\right]  :=\sum_{n\in\mathbb{N}}f_{n}a^{n}.
\]
(This sum is essentially finite, since $f$ is a polynomial.)

The element $f\left[  a\right]  $ is also denoted by $f\circ a$ or by
$f\left(  a\right)  $, and is called the \emph{value} of $f$ at $a$ (or the
\emph{evaluation} of $f$ at $a$, or the \emph{result of substituting }$a$ for
$x$ in $f$).
\end{definition}

\begin{lemma}
\label{lem.pol.eval.def}
\lean{FPS.eval_def}
\leanhelper
\leanok
The evaluation $f\left[a\right]$ can be computed as
$f\left[a\right]=\sum_{n}f_{n}\cdot a^{n}$,
where the sum ranges over all $n$ in the support of $f$.
\end{lemma}

\begin{proof}
\leanok
Immediate from the definitions.
\end{proof}

\begin{lemma}
\label{lem.pol.eval.finsum}
\lean{FPS.eval_eq_finsum}
\leanhelper
\leanok
The evaluation $f\left[a\right]$ equals
$\sum_{n\in\operatorname{supp}(f)}f_{n}\cdot a^{n}$,
making the finiteness of the sum explicit.
\end{lemma}

\begin{proof}
\leanok
This follows from the definition by restricting the sum to the support.
\end{proof}

\begin{lemma}
\label{lem.pol.eval.sum-range}
\lean{FPS.eval_eq_sum_range}
\leanhelper
\leanok
The evaluation $f\left[a\right]$ equals
$\sum_{n=0}^{\deg f}f_{n}\cdot a^{n}$,
where the sum is taken up to the degree of $f$.
\end{lemma}

\begin{proof}
\leanok
Since all coefficients beyond the degree of $f$ are zero, extending or
restricting the sum to this range does not change the result.
\end{proof}

\begin{theorem}
\label{thm.pol.eval.a+b}
\lean{FPS.eval_add', FPS.eval_mul', FPS.eval_smul', FPS.eval_C', FPS.eval_X', FPS.eval_X_pow', FPS.eval_comp'}
\leantarget
\leanok
Let $A$ be a $K$-algebra. Let $a\in A$. Then: \medskip

\textbf{(a)} Any $f,g\in K\left[  x\right]  $ satisfy%
\[
\left(  f+g\right)  \left[  a\right]  =f\left[  a\right]  +g\left[  a\right]
\ \ \ \ \ \ \ \ \ \ \text{and}\ \ \ \ \ \ \ \ \ \ \left(  fg\right)  \left[
a\right]  =f\left[  a\right]  \cdot g\left[  a\right]  .
\]


\textbf{(b)} Any $\lambda\in K$ and $f\in K\left[  x\right]  $ satisfy%
\[
\left(  \lambda f\right)  \left[  a\right]  =\lambda\cdot f\left[  a\right]
.
\]


\textbf{(c)} Any $\lambda\in K$ satisfies $\underline{\lambda}\left[
a\right]  =\lambda\cdot1_{A}$, where $1_{A}$ is the unity of the ring $A$.
\medskip

\textbf{(d)} We have $x\left[  a\right]  =a$. \medskip

\textbf{(e)} We have $x^{i}\left[  a\right]  =a^{i}$ for each $i\in\mathbb{N}%
$. \medskip

\textbf{(f)} Any $f,g\in K\left[  x\right]  $ satisfy $f\left[  g\left[
a\right]  \right]  =\left(  f\left[  g\right]  \right)  \left[  a\right]  $.
\end{theorem}

\begin{proof}
See \cite[Theorem 7.6.3]{19s} for parts \textbf{(a)}, \textbf{(b)},
\textbf{(c)}, \textbf{(d)} and \textbf{(e)}. Part \textbf{(f)} is
\cite[Proposition 7.6.14]{19s}.
\end{proof}

\subsection{Special evaluation results}

\begin{lemma}
\label{lem.pol.eval.X-self}
\lean{FPS.eval_X_eq_self}
\leanhelper
\leanok
For any polynomial $f\in K[x]$, we have $f[x]=f$.
\end{lemma}

\begin{proof}
\leanok
The evaluation map at $x$ is the identity on $K[x]$.
\end{proof}

\begin{lemma}
\label{lem.pol.eval.zero}
\lean{FPS.eval_zero_eq_coeff_zero}
\leanhelper
\leanok
For any polynomial $f\in K[x]$ and any $K$-algebra $A$, we have
$f[0]=\left[x^{0}\right]f$ (the constant term of $f$, coerced into $A$).
\end{lemma}

\begin{proof}
\leanok
Setting $a=0$: all terms $f_{n}\cdot 0^{n}$ for $n\geq 1$ vanish, leaving
only $f_{0}\cdot 0^{0}=f_{0}\cdot 1=f_{0}$.
\end{proof}

\begin{lemma}
\label{lem.pol.eval.one}
\lean{FPS.eval_one_eq_sum_coeffs}
\leanhelper
\leanok
For any polynomial $f\in K[x]$ and any $K$-algebra $A$, we have
$f[1]=\left[x^{0}\right]f+\left[x^{1}\right]f+\left[x^{2}\right]f+\cdots$
(the sum of all coefficients of $f$, coerced into $A$).
\end{lemma}

\begin{proof}
\leanok
Setting $a=1$: each term $f_{n}\cdot 1^{n}=f_{n}$, so
$f[1]=\sum_{n}f_{n}$.
\end{proof}
